\chapter{RANDOM WALKS}

\textit{July 22nd.}

Quizzes will account for 20\% of the final grade, and assignments for 5\%. The mid-term will account for 25\% of the final grade, and the final exam will account for 50\%. The primary reference for random walks is H.~Follner.

\section{The Setup}
We wish to concretify the idea of a random walk on the integers $\Z$. $\Omega$ is out sample space, with $0 < \abs{\Omega} < \infty$ (we shall later look at the case when $\Omega$ is countably infinite). $\cF$ is the event space, and is exactly $\cP(\Omega)$, the power set of $\Omega$. The intuition of a random walk is as follows; stand at the origin. At each step, flip a fair coin. If the coin comes up heads, take a step to the right; if tails, take a step to the left. Repeat this process $N$ times. $S_{n}$ here then denotes the position of the random walk after $n$ steps, for $0 \leq n \leq N$.

For a fixed $N \in \N$, we let $\Omega_{N} = \{\pm 1\}^{N}$, and $\cF_{N} = \cP(\Omega_{N})$. For $1 \leq k \leq N$, define the random variable $X_{k}:\Omega_{N} \to \{\pm 1\}$ as $X_{k}(\omega) = \omega_{k}$, where $\omega = (\omega_{1}, \omega_{2}, \ldots, \omega_{N}) \in \Omega_{N}$; this models the step taken by the random walk at step $k$. Further, define $S_{0} = 0$, and for $1 \leq n \leq N$, define $S_{n}:\Omega_{N} \to \Z$ as $S_{n}(\omega) = \sum_{k=1}^{n} X_{k}(\omega)$; this models the position of the random walk after $n$ steps. Finally, define the probability measure $\Prob_{N}$ on $(\Omega_{N}, \cF_{N})$ as $\Prob_{N}(A) = \frac{\abs{A}}{2^{N}}$ for $A \in \cF_{N}$. This models the fair coin flip at each step of the random walk.


\begin{definition}
    The sequence of random variables $\{S_{n}\}_{n=0}^{N}$ on $(\Omega_{N}, \cF_{N}, \Prob_{N} \equiv \Prob)$, as prescribed above, is called the \eax{simple symmetric random walk} of length $N$ on $\Z$, starting at the origin.
\end{definition}

Given $1 \leq k_{1} < k_{2} < \cdots < k_{n} \leq N$, and $x_{i} \in \{\pm 1\}$,
\begin{align}
    \Prob(X_{k_{1}} = x_{1}, \ldots, X_{k_{n}} = x_{n}) = \frac{\abs{\{\omega \in \Omega_{N} \mid X_{k_{i}}(\omega) = x_{i}, 1 \leq i \leq n\}}}{2^{N}} = \frac{2^{N-n}}{2^{N}} = 2^{-n}.
\end{align}
That is, fixing $n$ steps of the walk, the probability of any particular sequence of steps is $2^{-n}$. Using this, we can show that for any $1 \leq k \leq N$, $\Prob(X_{k} = 1) = \Prob(X_{k} = -1) = \frac{1}{2}$. Moreover, we can extract that
\begin{align}
    \Prob(X_{k}=x_{k},X_{l}=x_{l}) = 2^{-2} = \Prob(X_{k}=x_{k})\Prob(X_{l}=x_{l}),
\end{align}
showing independence of subsequent steps; the model seems like a reasonable fit for our intuition of a random walk.

\begin{example}
    The simple symmetric random walk has independent increments; that is, given $0 < k_{1} < k_{2} < N$, the random variables $S_{k}-S_{0}$, $S_{k_{2}}-S_{k_{1}}$, and $S_{N}-S_{k_{2}}$ are mutually independent. This is because the increments are functions of disjoint sets of independent random variables $X_{k}$, and hence are independent themselves. Similarly, one may concretely show that for $1 \leq k < m \leq N$, the distribution of $S_{m}-S_{k}$ is the same as that of $S_{m-k}$.
\end{example}

The following properties of the simple symmetric random walk are immediate consequences of the above definition and example.

Suppose we have $\Prob(S_{n-1}=a_{n-1}, S_{n-2}=a_{n-2}, \ldots, S_{0}=0) > 0$. Then, we are free to compute the conditional probability $\Prob(S_{n}=a_{n} \mid S_{n-1}=a_{n-1}, S_{n-2}=a_{n-2}, \ldots, S_{0}=0)$. One may verify that
\begin{align}
    \Prob(S_{n}=a_{n} \mid S_{n-1}=a_{n-1}, S_{n-2}=a_{n-2}, \ldots, S_{0}=0) = \Prob(S_{n}=a_{n} \mid S_{n-1}=a_{n-1}),
\end{align}
for all $a_{1},\ldots,a_{n} \in \Z$ such that $\Prob(S_{n-1}=a_{n-1}, S_{n-2}=a_{n-2}, \ldots, S_{0}=0) > 0$. This is called the \eax{Markov property} of the simple symmetric random walk. In words, this property states that given the present state of the walk, the future evolution of the walk is independent of its past history.

Now suppose $\Prob(S_{k}=a) > 0$ for some $k \geq 1,a \in \Z$, and let $m > k,b \in \Z$. Then, one can simply compute that $\Prob(S_{m}=b \mid S_{k}=a) = \Prob(S_{m-k}=b-a)$.

The next question we may ask is, on an average, where the random walk is located after $n$ steps, for $1 \leq n \leq N$. We can compute the expected value of $S_{n}$ as follows:
\begin{align}
    \E[S_{n}] = \E\left[\sum_{k=1}^{n} X_{k}+0\right] = \sum_{k=1}^{n} \E[X_{k}] = \sum_{k=1}^{n} 0 = 0. 
\end{align}
Thus, to get a better idea of the location of the random walk after $n$ steps, we compute the variance of $S_{n}$:
\begin{align}
    \Var(S_{n}) = \Var\left(\sum_{k=1}^{n} X_{k}+0\right) = \sum_{k=1}^{n} \Var(X_{k}) = \sum_{k=1}^{n} 1 = n.
\end{align}
So, via the central limit theorem, we can approximate for large $n$ as
\begin{align}
    \Prob\left( -3 \leq \frac{S_{n}}{\sqrt{n}} \leq 3 \right) \approx \Phi(3) - \Phi(-3)
\end{align}
where $\Phi$ is the cumulative distribution function of the standard normal distribution.

We now claim that
\begin{align}
    \Prob(S_{2n}=0) = \binom{2n}{n} 2^{-2n} = \Prob(S_{2n-1}=1) = \Prob(S_{2n-1}=-1).
\end{align}
To this end, we note that for $S_{n} = x$, we must have taken $\frac{n+x}{2}$ steps to the right, and $\frac{n-x}{2}$ steps to the left. Thus, we have
\begin{align}
    \Prob(S_{n}=x) = \frac{\abs{\{\omega \in \Omega_{N} \mid S_{n} = x\}}}{2^{N}} = \binom{n}{\frac{n+x}{2}} \frac{2^{N-n}}{2^{N}} = \binom{n}{\frac{n+x}{2}} 2^{-n}.
\end{align}
The binomial coefficient attains its maximum at $x=0$ for even $n$, and at $x=\pm 1$ for odd $n$. Thus, we have shown the above required claim. We can further approximate this via \eax{Stirling's formula}, which stats $k! \sim \sqrt{2\pi k} k^{k}e^{-k}$ for large $k$. For now, we shall treat this as an equality. Thus, for large $n$,
\begin{align}
    \Prob(S_{2n}=0) = \binom{2n}{n} 2^{-2n} \sim \frac{\sqrt{4\pi n}(2n)^{2n}e^{-2n}}{2\pi n n^{2n} e^{-2n}} 2^{-2n} = \frac{1}{\sqrt{\pi n}}.
\end{align}
We now show that given any interval $[a,b] \subsetneq \R$, the probability that the random walk stays in this interval is zero as $n \to \infty$. To this end, note that either $\Prob(S_{n}=x) \leq \Prob(S_{n}=0)$ or $\Prob(S_{n}=x) \leq \Prob(S_{n}=1)$ holds for any given $n$; moreover, both of the quantities on the right-hand side are bounded by $\frac{c}{\sqrt{n}}$ for an appropriate constant $c > 0$. Thus, we have
\begin{align}
    0 \leq \Prob(a \leq S_{n} \leq b) = \sum_{x=a}^{b} \Prob(S_{n}=x) \leq \frac{c(b-a+1)}{\sqrt{n}} \to 0 \text{ as } n \to \infty.
\end{align}


\section{Stopping Times}
\noindent \textit{July 29th.}

We can interpret $S_{n}$ as the amount of capital earned by a player until time $n$. Earlier, we saw that the expected `amount of capital' at $n$ is $\E[S_{n}] = 0$, for $0 \leq n \leq N$. We ask if there is a favourite (random) time at which the player should stop playing, so as to maximize the expected amount of capital earned. We assume that there is no insider trading, and that the player's decision to stop playing at time $n$ is based on the information available until time $n$.

\begin{definition}
    A subset $A \subseteq \Omega$ is termed an \eax{observable event} by time $n$ if it can be written as a union of basic events of the form
    \begin{align}
        \{\omega \in \Omega \mid \omega_{1}=o_{1}, \ldots, \omega_{n}=o_{n}\}, \quad o_{1},\ldots,o_{n} \in \{\pm 1\}.
    \end{align}
    That is, $A$ can be determined by observing the outcomes of the first $n$ coin tosses.
\end{definition}

We shall use $\cA_{n}$ to denote the collection of all observable events by time $n$. One can show that
\begin{align}
    \{\emptyset,\Omega\} = \cA_{0} \subseteq \cA_{1} \subseteq \cdots \subseteq \cA_{N} = \cF.
\end{align}

\begin{definition}
    A map $T:\Omega \to \{0,1,\ldots,N\} \cup \{\infty\}$ is called a \eax{stopping time} if $\{T=n\} \in \cA_{n}$ for all $0 \leq n \leq N$. 
\end{definition}

As an example, pick $a \in \Z$. We can then define $\sigma_{a}:\Omega \to \{0,1,2,\ldots,N\} \cup \{\infty\}$ as $\sigma_{a}(\omega) = \min\{n \geq 0 \mid S_{n}(\omega) = a\}$, with the convention that $\min \emptyset = \infty$. This is, simply put, the first time the random walk hits the point $a$. The reader may show that $\sigma_{a}$ is indeed a stopping time. On the other hand, the map $\tau:\Omega \to \{0,1,2,\ldots,N\} \cup \{\infty\}$ defined as $\tau(\omega) = \min\{n \geq 0 \mid S_{n}(\omega) = S_{N}(\omega)\}$ is not a stopping time; this is because we cannot determine if the random walk has hit its final position at time $n$ without knowing the outcome of the coin tosses after time $n$.

\begin{theorem}
    For any stopping time $T:\Omega \to \{0,1,\ldots,N\}$, we have $\E[S_{T}] = 0$.
\end{theorem}
That is, there is an impossibility of a favourite stopping time. We shall adopt the notation $S_{T}(\omega) \defeq S_{T(\omega)}(\omega)$, the outcome of the trajectory $\omega$ at the stopping time $T(\omega)$. This leads us to discuss the impossibillity of profitable games.


Let $V_{1},\ldots,V_{N}:\Omega \to \R$ be random variables such that $\{V_{k}=c\} \in \cA_{k-1}$ for all $c \in \R$, $k = 1,\ldots,N$. $V_{k}$ here represents the amount of capital a player bets at time $k$. The result at the $k$-th round of the game is then $V_{k}X_{k}$. This is known as a game system.

\begin{theorem}
    For any $V_{1},\ldots,V_{N}$ which are a game system as above, let $S_{n}^{V} = \sum_{k=1}^{n} V_{k}X_{k}$ for $1 \leq n \leq N$. Then, $\E[S_{N}^{V}] = 0$.
\end{theorem}

\begin{proof}
    We have $\E[S_{N}^{V}] = \sum_{k=1}^{N} \E[V_{k}X_{k}]$. Let us restrict our range as $\mathrm{range}(V_{k}) = \{c_{1},\ldots,c_{M}\}$ for some $M \in \N$. Then, we have
    \begin{align}
        \E[V_{k}X_{k}] = \sum_{i=1}^{M} c_{i} \, \E[\1_{\{V_{k}=c_{i}\}} X_{k}]
    \end{align}
    Note that, by definition, $\{V_{k}=c_{i}\} \in \cA_{k-1}$, and hence is independent of $X_{k}$. Thus, we have
    \begin{align}
        \E[V_{k}X_{k}] = \sum_{i=1}^{M} c_{i} \, \E[\1_{\{V_{k}=c_{i}\}}] \, \E[X_{k}] = 0.
    \end{align}
\end{proof}

We shall use this theorem to prove the theorem before it.

\begin{proof}[Proof of impossibility of a favourite stopping time]
    We have
    \begin{align}
        S_{T} = \sum_{k=1}^{T(\omega)} X_{k} = \sum_{k=1}^{N} \1_{\{T \geq k\}} X_{k}.
    \end{align}
    Now, for $1 \leq k \leq N$, set $V_{k}:\Omega \to \{0,1\}$ as $V_{k}(\omega) = \1_{\{T(\omega) \geq k\}}$. Then $\{V_{k} = 1\} = \{T \geq k\} = \{T < k\}^{c} = \left( \bigcup_{m=0}^{k-1} \{T=m\} \right)^{c} \in \cA_{k-1}$ since $\{T=m\} \in \cA_{m} \subseteq \cA_{k-1}$ for $0 \leq m \leq k-1$, and $cA_{k-1}$ is closed under finite unions and complements. Thus, $V_{1},\ldots,V_{N}$ is a game system, and hence $\E[S_{T}] = \E[S_{N}^{V}] = 0$ by the previous theorem.
\end{proof}

Suppose $T:\Omega \to \{0,1,\ldots,N\}$ is a stopping time. Then $\Var(S_{T}) = \E[S_{T}^{2}] - \E[S_{T}]^{2} = \E[S_{T}^{2}]$. We can compute $\E[S_{T}^{2}]$ as follows; let $1 \leq k \leq N$. Then
\begin{align}
    S_{k}^{2} = (S_{k-1}+X_{k})^{2} = S_{k-1}^{2} + 2X_{k}S_{k-1} + 1.
\end{align}
Set $V_{k} = S_{k-1}\1_{\{T \geq k\}}$. Verify that $V_{1},\ldots,V_{N}$ is a game system. Then, we have
\begin{align}
    S_{N}^{V} &= \sum_{k=1}^{N} X_{k}V_{k} = \sum_{k=1}^{N} (S_{k-1}X_{k}) \1_{\{T \geq k\}} = \sum_{k=1}^{N} \frac{1}{2}(S_{k}^{2} - S_{k-1}^{2}-1)\1_{\{T \geq k\}} \notag \\
    &= \frac{1}{2}\sum_{k=1}^{N} S_{k}^{2} \1_{\{T \geq k\}} - \frac{1}{2} \sum_{k=1}^{N} S_{k-1}^{2} \1_{\{T \geq k\}} - \frac{1}{2}\sum_{k=1}^{N} \1_{\{T \geq k\}} \notag \\
    &= \frac{1}{2}\sum_{k=1}^{N} S_{k}^{2} \1_{\{T=k\}} - \frac{T}{2} = \frac{S_{T}^{2}-T}{2}.
\end{align}
By the previous theorem, we have $\E[S_{N}^{V}] = 0$, and hence $\E[S_{T}^{2}] = \E[T]$. 

\subsection{Ruin Problem}
We have two players, one with capital $a$ and the other with capital $b$. For $0 \leq n \leq N$, let $S_{n}$ denote the gain of the first player. We define the first player's ruin as
\begin{align}
    \{\sigma_{-a} < \sigma_{b},\, \sigma_{-a} \leq N\}.
\end{align}
Define the ruin probability $\gamma_{N}^{1}$ to be the probability of this event, the first player's ruin. Similarly, we define the second player's ruin as
\begin{align}
    \{\sigma_{b} < \sigma_{-a},\, \sigma_{b} \leq N\},
\end{align}
and define the ruin probability $\gamma_{N}^{2}$ to be the probability of this event, the second player's ruin. Now define $T_{N}:\Omega \to \{0,1,2,\ldots,N\}$ as $T_{N} = \min\{\sigma_{-a},\sigma_{b},N\}$. Then, $T_{N}$ is a stopping time. It has been shown that $\E[S_{T_{N}}] = 0$. Explicitly writing out $S_{T_{N}}$, we have
\begin{align}
    S_{T_{N}} = \begin{cases}
        -a &\text{if } \sigma_{-a} < \sigma_{b} \text{ and } \sigma_{-a} \leq N, \\
        b &\text{if } \sigma_{b} < \sigma_{-a} \text{ and } \sigma_{b} \leq N, \\
        S_{N} &\text{if } \min(\sigma_{-a},\sigma_{b}) \geq N.
    \end{cases}
\end{align}
Taking expectations, we have
\begin{align}
    0 &= -a\Prob(\sigma_{-a} < \sigma_{b},\, \sigma_{-a} \leq N) + b\Prob(\sigma_{b} < \sigma_{-a},\, \sigma_{b} \leq N) + \E[S_{N}\1_{\{\min(\sigma_{-a},\sigma_{b}) \geq N\}}] \notag \\
    \implies 0 &= -a\gamma_{N}^{1} + b\gamma_{N}^{2} + \E[S_{N}\1_{\{\min(\sigma_{-a},\sigma_{b}) \geq N\}}].
\end{align}
Here, we claim that $\E[S_{N}\1_{\{\min(\sigma_{-a},\sigma_{b}) \geq N\}}] \to 0$ as $N \to \infty$. To this end, first notice the bounds
\begin{align}
    0 \leq \abs{S_{n}} \1_{\{\min(\sigma_{-a},\sigma_{b}) \geq N\}} \leq (b+a) \1_{\{\min(\sigma_{-a},\sigma_{b}) \geq N\}}
\end{align}
which then implies
\begin{align}
    0 \leq \E[\abs{S_{n}} \1_{\{\min(\sigma_{-a},\sigma_{b}) \geq N\}}] \leq (b+a) \Prob(\min(\sigma_{-a},\sigma_{b}) \geq N) \leq (b+a) c\frac{(b-a+1)}{\sqrt{N}} \to 0 \text{ as } N \to \infty.
\end{align}
We now make three `leaps of faith;' they are
\begin{enumerate}
    \item $\gamma^{1} = \lim_{N \to \infty} \gamma_{N}^{1}$ and $\gamma^{2} = \lim_{N \to \infty} \gamma_{N}^{2}$ exist.
    \item $\gamma^{1} + \gamma^{2} = 1$.
    \item $0 = -a\gamma^{1} + b\gamma^{2}$.
\end{enumerate}
The last two assumptions imply that $\gamma^{1} = \frac{b}{a+b}$ and $\gamma^{2} = \frac{a}{a+b}$.

\section{Reflection Principle}

Consider a random walk starting at the origin that ends up at $a+c$. Such a walk will have to pass through $a$ at some point. We can then reflect the part of the walk after it hits $a$ about the line $y=a$. This gives us a bijection between walks that end up at $a+c$ and walks that reach $a$ and then end up at $a-c$. That is,
\begin{align}
    \abs{\{\omega \in \Omega \mid S_{n} = a+c\}} = \abs{\{\omega \in \Omega \mid S_{n} = a+c,\, \sigma_{a} \leq n\}} = \abs{\{\omega \in \Omega \mid S_{n} = a-c,\, \sigma_{a} \leq n\}}.
\end{align}

This lemma immediately follows.
\begin{lemma}
    $\Prob(S_{n}=a+c) = \Prob(\sigma_{a} \leq n,\, S_{n}=a-c)$.
\end{lemma}

\begin{theorem}
    $\Prob(\sigma_{a} \leq n) = \Prob(S_{n} \notin [-a,a-1])$, and $\Prob(\sigma_{a}=n) = \frac{a}{n}\Prob(S_{n}=a)$.
\end{theorem}
\begin{proof}
    We have
    \begin{align}
        \Prob(\sigma_{a} \leq n) &= \sum_{b \in \Z} \Prob(\sigma_{a} \leq n,\, S_{n} = b) = \sum_{b \in \Z,\, b \geq a} \Prob(\sigma_{a} \leq n,\, S_{n} = b) + \sum_{b \in \Z,\, b < a} \Prob(\sigma_{a} \leq n,\, S_{n} = b) \notag \\
        &= \Prob(S_{n} \geq a) + \sum_{b \in \Z,\, b < a} \Prob(S_{n}=2a-b) = \Prob(S_{n} \geq a) + \Prob(S_{n} > a) = \Prob(S_{n} \geq a) + \Prob(S_{n} < -a) \notag \\
        &= \Prob(S_{n} \notin [-a,a-1]).
    \end{align}
    On the other hand,
    \begin{align}
        \Prob(\sigma_{a}=n) &= \Prob(\sigma_{a} \leq n) - \Prob(\sigma_{a} \leq n-1) = \frac{1}{2} \left( \Prob(S_{n-1} = a-1) - \Prob(S_{n-1}=a+1)\right) \notag \\
        &= \frac{1}{2} \left( \binom{n-1}{\frac{n-1+a-1}{2}}\frac{1}{2^{n-1}} - \binom{n-1}{\frac{n-1-a+1}{2}}\frac{1}{2^{n-1}} \right) = \frac{1}{2^{n}} \cdot \frac{a}{n} \binom{n}{\frac{n+a}{2}} = \frac{a}{n}\Prob(S_{n}=a).
    \end{align}
\end{proof}

We now take as convention that $\sigma_{0} = \min\{n \geq 1 \mid S_{n}=0\}$ so that $\sigma_{0} \neq 0$.

\begin{lemma}
    $\Prob(\sigma_{0}>2n) = \Prob(S_{2n-1}=-1)$.
\end{lemma}
\begin{proof}
    We have
    \begin{align}
        \Prob(\sigma_{0} > 2n) &= \Prob(S_{1} \neq 0, S_{2} \neq 0, \ldots, S_{2n} \neq 0) = 2\Prob(S_{1} > 0, \ldots, S_{2n} > 0) \notag \\
        &= 2 \cdot \frac{1}{2^{2n}} \#\{\text{paths that start at $0$ and don't hit $-1$ in $2n-1$ steps}\} \notag \\
        &= \frac{1}{2^{2n-1}} \#\{\text{paths that start at $0$ and stay below $1$ in $2n-1$ steps}\} = \Prob(\sigma_{1} > 2n-1) \notag \\
        &= 1-\Prob(\sigma_{1} \leq 2n-1) = 1-\Prob(S_{2n-1} \notin [0,1]) = 1-\Prob(S_{2n-1} \neq -1) = \Prob(S_{2n-1}=-1).
    \end{align}
\end{proof}

\subsection{Arcsine Law}
\textit{August 5th.}

The arcsine law is concerned with the analysis of the last visit to the origin till $2N$. That is, defining $L = \max\{0 \leq n \leq 2N \mid S_{n}=0\}$, we wish to compute the order of $L$.

\begin{theorem}
    Let $n \in \N_{0}$, with $n \leq N$. Then $\Prob(L = 2n) = \frac{1}{2^{2N}} \binom{2n}{n} \binom{2N-2n}{N-n}$.
\end{theorem}

\begin{remark}
    Suppose this theorem holds; using Stirling's formula, that $n! \sim n^{n}e^{-n}$ for large $n$, we can approximate $\Prob(L = 2n) \sim \frac{1}{\pi \sqrt{n(N-n)}} = \frac{1}{\pi N} \frac{1}{\sqrt{\frac{n}{N}(1-\frac{n}{N})}}$. The idea is to find the order of $\Prob(L \leq 2Nx) = \Prob(\frac{L}{2N} \leq x)$ for $x \in [0,1]$. Rewriting gives
    \begin{align}
        \Prob(L \leq 2Nx) = \sum_{n \leq \floor{Nx}} \Prob(L = 2n) = \sum_{n \leq \floor{Nx}} \frac{1}{N} f(\frac{1}{N})
    \end{align}
    where $f(x) = \frac{1}{\pi \sqrt{x(1-x)}}$. This is a Riemann sum, and hence we have $\Prob(L \leq 2Nx) \sim \int_{0}^{x} f(t) dt = \frac{2}{\pi} \arcsin(\sqrt{x})$.
\end{remark}

\begin{proof}[Proof of theorem]
    We simply have
    \begin{align}
        \Prob(L = 2n) = \Prob(S_{2n}=0)\Prob(\sigma_{0} > 2N-2n) = \binom{2n}{n} \frac{1}{2^{2n}} \cdot \Prob(S_{2N-2n}=0) = \frac{1}{2^{2N}} \binom{2n}{n} \binom{2N-2n}{N-n}.
    \end{align}
\end{proof}

\section{Higher Dimensions and Infinite Walks}
We are now fit to make sense of simple random walks on higher dimensions; let $\Z^{d} = \{(x_{1},\ldots,x_{d}) \mid x_{i} \in \Z, 1 \leq i \leq d\}$ for $d \geq 1$. Also, let $\abs{x} = \sqrt{\sum_{i=1}^{d} x_{i}^{2}}$. Then, for a fixed $N \geq 1$, we define
\begin{align}
    \Omega_{N} = \{\omega = (\omega_{1},\ldots,\omega_{N}) \mid \omega_{i} \in \Z^{d}, \, \abs{\omega_{i}} = 1,\, 1 \leq i \leq N\}
\end{align}
and $\xi_{i}:\Omega_{N} \to \Z^{d}$ as $\xi_{i}(\omega) = \omega_{i}$ for $1 \leq i \leq N$. We can then further define $X_{n}:\Omega_{N} \to \Z^{d}$ as $X_{n}(\omega) = \sum_{i=1}^{n} \xi_{i}(\omega)$ for $1 \leq n \leq N$. Moreover, $\cF = \cP(\Omega_{N})$, and $\Prob(A) = \frac{\abs{A}}{(2d)^{N}}$ for $A \in \cF$. This is the simple symmetric random walk on $\Z^{d}$ of length $N$.

We wish to extend our notion of random walks for ones of infinite length; that is, the consistency of $\Prob_{N}$ as $N$ varies. We define the projection map $\Pi_{M}:\Omega_{N} \to \Omega_{M}$ as $\omega \mapsto (\omega_{1},\ldots,\omega_{M}) =\vcentcolon \tilde{\omega}$. Now fix $\tilde{\omega} \in \Omega_{M}$ so that
\begin{align}
    \Prob_{N}(\{\omega \in \Omega_{N} \mid \Pi_{M}(\omega) = \tilde{\omega}\}) = \frac{\abs{\{\omega \in \Omega_{N} \mid \Pi_{N}(\omega) = \tilde{\omega}\}}}{(2d)^{N}} = \frac{(2d)^{N-M}}{(2d)^{N}} = (2d)^{-M} = \Prob_{M}(\{\tilde{\omega}\}).
\end{align}

\begin{theorem}[\eax{Kolmogorov consistency theorem}]
    There exists a probability space $(\Omega,\cF,\Prob)$ such that $\Omega \equiv \Omega_{\infty} = \{(\omega_{i})_{i \geq 1} \mid \omega_{i} \in \Z^{d}, \, \abs{w_{i}} = 1\}$, with $\cF$ suitabily defined, and $\Prob:\cF \to [0,1]$ such that $\Prob(\omega \in \Omega \mid \Pi_{M}(\omega) = \tilde{\omega}) = \Prob_{M}(\{\tilde{\omega}\})$ for all $M \geq 1$ and $\tilde{\omega} \in \Omega_{M}$.
\end{theorem}

A proof of this theorem is beyond the scope of this course; this lays the groundwork for the study of random walks of infinite length, by defining $\xi_{i}:\Omega \to \Z^{d}$ as $\xi_{i}(\omega) = \omega_{i}$ for $i \geq 1$, and $X_{n}:\Omega \to \Z^{d}$ as $X_{n}(\omega) = \sum_{i=1}^{n} \xi_{i}(\omega)$ for $n \geq 1$. This is the simple symmetric random walk on $\Z^{d}$ of infinite length.

\subsection{Large Number Results}
Let $\{X_{n}\}_{n \geq 1}$ be the simple symmetric random walk on $\Z^{d}$ of infinite length. Define
\begin{align}
    A = \{\omega \in \Omega \mid X_{n}(\omega)/n \to 0 \text{ as } n \to \infty\}.
\end{align}
Then, one can show that $\Prob(A) = 1$ despite $A \neq \Omega$. This is a consequence of \eax{Stirling's law of large numbers}. Also,
\begin{align}
    \Prob\left( \frac{X_{n}}{\sqrt{n}} \leq x \right) \to \int_{-\infty}^{x} \frac{e^{-y^{2}/2}}{\sqrt{2\pi}} dy \text{ as } n \to \infty.
\end{align}
for when $d = 1$. This is a consequence of the \eax{central limit theorem}.